\documentclass[10pt,a4paper]{article}

\usepackage[utf8]{inputenc}
\usepackage[T1]{fontenc}
\usepackage[brazilian]{babel}
\usepackage[a4paper,top=1.25cm,bottom=1.35cm,left=1.5cm,right=1.5cm,headheight=13pt]{geometry}
\usepackage{xcolor}
\usepackage{enumitem}
\usepackage{fancyhdr}
\usepackage{hyperref}
\usepackage[most]{tcolorbox}

\setlength{\emergencystretch}{2em}

\definecolor{BrandDark}{HTML}{1E293B}
\definecolor{BrandSlate}{HTML}{475569}
\definecolor{BrandAmber}{HTML}{B45309}
\definecolor{BrandBg}{HTML}{F8FAFC}
\definecolor{BrandLine}{HTML}{CBD5E1}
\definecolor{OkGreen}{HTML}{15803D}
\definecolor{InfoBlue}{HTML}{1D4ED8}

\hypersetup{
  colorlinks=true,
  linkcolor=InfoBlue,
  urlcolor=InfoBlue,
  pdftitle={Roteiro Explicativo — Integrante 1 (Slides 1 a 3)},
  pdfauthor={Grupo BD2 (UFPI/CSHNB)}
}

\setlength{\parindent}{0pt}
\setlength{\parskip}{3pt}

\pagestyle{fancy}
\fancyhf{}
\fancyhead[L]{\footnotesize\textcolor{BrandSlate}{\textbf{BD II (2026.2)} · Guia Individual do Apresentador}}
\fancyhead[R]{\footnotesize\textcolor{BrandSlate}{\textbf{Integrante 1} (Slides 1--3)}}
\fancyfoot[L]{\footnotesize\textcolor{BrandSlate}{Orçamento humano alvo: $\sim$2m00s (teto máximo: 2m30s)}}
\fancyfoot[R]{\footnotesize\textcolor{BrandSlate}{Pág. \thepage}}
\renewcommand{\headrulewidth}{0.4pt}
\renewcommand{\footrulewidth}{0.4pt}

\newcommand{\tagTime}[1]{\colorbox{BrandDark}{\textcolor{white}{\textbf{\footnotesize\ #1\ }}}}

\newcommand{\slideHeader}[3]{%
  \noindent{\Large\bfseries\color{BrandDark} Slide #1 — #2}\hfill\tagTime{Alvo: #3}\par
  \vspace{2pt}\noindent\textcolor{BrandAmber}{\rule{\linewidth}{1.3pt}}\par
  \vspace{4pt}
}

\newtcolorbox{telaBox}[1][O que aparece neste slide no Canva]{
  colback=BrandBg,
  colframe=BrandSlate!60!black,
  fonttitle=\bfseries\small,
  title={#1},
  boxrule=0.6pt,
  arc=3pt,
  left=8pt,right=8pt,top=5pt,bottom=5pt
}

\newtcolorbox{falaBox}[1][Fala sugerida (leitura natural em voz alta)]{
  colback=InfoBlue!4!white,
  colframe=InfoBlue!75!black,
  fonttitle=\bfseries\small,
  title={#1},
  boxrule=0.8pt,
  arc=3pt,
  left=9pt,right=9pt,top=6pt,bottom=6pt
}

\newtcolorbox{expBox}[1][Explicação em tópicos das informações da fala (cola técnica para o apresentador)]{
  colback=OkGreen!4!white,
  colframe=OkGreen!75!black,
  fonttitle=\bfseries\small,
  title={#1},
  boxrule=0.7pt,
  arc=3pt,
  left=8pt,right=8pt,top=6pt,bottom=6pt
}

\begin{document}

% Cabeçalho do Integrante 1
\begin{tcolorbox}[colback=BrandDark,colframe=BrandDark,arc=4pt,left=11pt,right=11pt,top=8pt,bottom=8pt]
  {\Large\bfseries\color{white} Roteiro Explicativo — Integrante 1 (Slides 1 a 3)}\\[3pt]
  {\small\color{BrandLine} Contexto Rápido, Problema de Pesquisa e Fundamentos de B+ Tree (SQL Server 2022)}\\[4pt]
  {\footnotesize\color{white}
    \textbf{Tempo alvo:} $\sim$2m00s \quad|\quad
    \textbf{Teto:} 2m30s \quad|\quad
    \textbf{Ritmo:} $\sim$120 palavras/min \quad|\quad
    \textbf{Próximo a falar:} Integrante 2 (Slide 4)
  }
\end{tcolorbox}

\vspace{4pt}
{\small\color{BrandSlate}\textbf{Regra desta folha:} cada página traz \textbf{um único slide por vez}, evitando confusão na hora da fala. Leia o bloco azul com calma e use o bloco verde caso a banca faça perguntas ou você queira explicar com suas próprias palavras.}

\vspace{8pt}

% ==============================================================================
% SLIDE 1
% ==============================================================================
\slideHeader{1}{Capa e Abertura}{$\sim$0m20s}

\begin{telaBox}
  \begin{itemize}[leftmargin=14pt,itemsep=1pt,topsep=1pt]
    \item \textbf{Título:} \emph{Indexação e Otimização no SQL Server 2022}.
    \item \textbf{Subtítulo:} Planos de execução e sargabilidade em 6,25M linhas $\cdot$ BD II — UFPI (CSHNB).
    \item \textbf{Grupo:} Hermeson A., João Marcos R., Lizzandro W., Marcos W. e Raildom R.
  \end{itemize}
\end{telaBox}

\begin{falaBox}
``Bom dia, professor e colegas. No nosso projeto final de Banco de Dados II, investigamos experimentalmente como escolhas reais de indexação, reescrita SQL e operadores físicos alteram o desempenho do Microsoft SQL Server 2022 em uma base de e-commerce com 6,25 milhões de registros. Nossa apresentação rápida de slides trará o contexto, como o experimento foi construído e os principais resultados, deixando metade do tempo total livre para a demonstração prática ao vivo. Eu abro com o problema e os fundamentos.''
\end{falaBox}

\begin{expBox}
  \begin{itemize}[leftmargin=14pt,itemsep=2pt,topsep=1pt]
    \item \textbf{Por que mencionar ``6,25 milhões de linhas'' logo na capa?} Para deixar claro desde o início que o estudo não usou uma tabelinha didática de dezenas de linhas — onde tudo caberia em 1 única página e índices quase não fariam diferença —, mas sim um volume grande o suficiente para expor custos reais de páginas em disco e memória.
    \item \textbf{Divisão dos 20 minutos totais da equipe:} até 10 min máximos para os slides teóricos (3 integrantes; orçado em $\sim$8 min para folga humana) e $\sim$9--10 min dedicados à prática ao vivo no VS Code + Docker (2 integrantes).
    \item \textbf{Disciplina e orientação:} Projeto final da disciplina Banco de Dados II (2026.2), Bacharelado em Sistemas de Informação, UFPI/CSHNB (Picos -- PI), Prof. Dennis Sávio.
  \end{itemize}
\end{expBox}

\clearpage

% ==============================================================================
% SLIDE 2
% ==============================================================================
\slideHeader{2}{Contexto e Problema}{$\sim$0m50s}

\begin{telaBox}
  \begin{itemize}[leftmargin=14pt,itemsep=1pt,topsep=1pt]
    \item 5 pontos: (1) Gargalo de E/S lógica $O(N)$ vs. \emph{buffer pool}; (2) Armadilha do \emph{Key Lookup}; (3) Predicados não \emph{sargable}; (4) Mito do ``índice gratuito''; (5) Objetivo no SQL Server 2022 (10 leituras + 1 escrita).
  \end{itemize}
\end{telaBox}

\begin{falaBox}
``Quando aplicações crescem para milhões de linhas, o principal gargalo não é apenas o disco lento: varreduras integrais de tabela, de complexidade $O(N)$, consomem páginas em memória, pressionam o \emph{buffer pool} e travam a escalabilidade mesmo em SSDs rápidos. Na prática, muita gente tenta resolver isso criando índices sem critério e cai em três armadilhas. Primeira: índices simples sem as colunas projetadas geram centenas de \emph{Key Lookups} aleatórios, fazendo o otimizador desistir e voltar para a varredura completa. Segunda: filtros com funções ou curinga no início impedem o uso do índice. E terceira: existe o mito do índice gratuito — todo índice extra precisa ser atualizado nas escritas. Nosso objetivo foi medir, com rigor estatístico no SQL Server 2022, esses ganhos e contra-exemplos em dez leituras e uma carga de escrita em lote.''
\end{falaBox}

\begin{expBox}
  \begin{itemize}[leftmargin=14pt,itemsep=2pt,topsep=1pt]
    \item \textbf{O que significa ``E/S lógica'' (\emph{logical reads})?} É o número de páginas de 8 KB que o SQL Server lê a partir do \emph{buffer pool} (memória RAM reservada ao cache de dados) para responder à consulta. Mesmo sem acessar o SSD físico (\emph{Warm Cache}), ler 8.090 páginas em vez de 5 páginas desperdiça memória, largura de banda interna e CPU.
    \item \textbf{Por que um SSD rápido não ``cura'' uma varredura $O(N)$?} Porque percorrer toda a árvore clusterizada ocupa páginas úteis do cache (expulsando outros dados quentes) e mantém a CPU varrendo centenas de milhares de tuplas inúteis linha a linha.
    \item \textbf{Qual a ``armadilha do Key Lookup''?} Se o índice guarda apenas a chave de busca, mas o \texttt{SELECT} pede outras colunas, o motor usa a árvore secundária só para achar a chave e depois pula de volta para a tabela principal para buscar cada linha completa. Com dezenas/centenas de linhas, esses pulos aleatórios ficam tão caros que o otimizador volta para o \emph{Clustered Index Scan}.
    \item \textbf{O que é o ``mito do índice gratuito''?} Em leitura (\texttt{SELECT}), o índice acelera a busca; mas em \texttt{INSERT/UPDATE/DELETE}, o SGBD precisa manter cada árvore B+ sincronizada na mesma transação (\emph{write amplification}), adicionando E/S, trava de log e possíveis divisões de página (\emph{page splits}).
  \end{itemize}
\end{expBox}

\clearpage

% ==============================================================================
% SLIDE 3
% ==============================================================================
\slideHeader{3}{Fundamentos: Árvore B+ e Planos de Execução}{$\sim$0m50s}

\begin{telaBox}
  \begin{itemize}[leftmargin=14pt,itemsep=1pt,topsep=1pt]
    \item 5 conceitos-chave: Árvore B+ (páginas de 8 KB, $O(\log_m N)$); Índice Clusterizado vs. Não-clusterizado (\texttt{INCLUDE}); Composto $(K_1, K_2)$ e Filtrado; Sargabilidade e \emph{Anti Semi Join}; \emph{Nested Loops}, \emph{Merge}, \emph{Hash Match} e \emph{Batch Mode}.
  \end{itemize}
\end{telaBox}

\begin{falaBox}
``Para entender nossos planos de execução, bastam cinco peças teóricas. No SQL Server, índices usam árvores B+ em páginas fixas de 8 kilobytes: os nós internos só roteiam com alto fator de ramificação, permitindo localizar um registro entre milhões de linhas em apenas 2 a 4 páginas, em $O(\log_m N)$. Cada tabela possui no máximo um índice clusterizado, cujas folhas armazenam as próprias linhas ordenadas; os demais são não-clusterizados, e podem incorporar colunas extras nas folhas via cláusula \texttt{INCLUDE} para formar índices de cobertura. Também usamos índices compostos, em que a igualdade vem antes da faixa de valores, e índices filtrados com cláusula \texttt{WHERE}. Já a \emph{sargabilidade} garante que a coluna filtrada fique isolada no predicado para viabilizar o \emph{Index Seek}. Por fim, além de \emph{Nested Loops}, \emph{Merge} e \emph{Hash Match}, o SQL Server 2022 possui o modo vetorizado \emph{Batch Mode on Rowstore}, que processa lotes de até 900 linhas — detalhe decisivo no caso Q7. Agora o Integrante 2 mostra como montamos o projeto na prática.''
\end{falaBox}

\begin{expBox}
  \begin{itemize}[leftmargin=14pt,itemsep=2pt,topsep=1pt]
    \item \textbf{Por que a Árvore B+ tem altura de apenas 2 a 4 níveis para milhões de linhas?} Porque diferentemente da Árvore B clássica, os nós internos da B+ não carregam os registros completos — guardam apenas chaves e ponteiros. Assim cabem centenas de ponteiros por página de 8 KB (alto \emph{fan-out} $m$): $\log_m(N)$ cresce muito devagar.
    \item \textbf{Diferença rápida — Clusterizado vs. Não-clusterizado:}
      \emph{Clusterizado} é a própria tabela organizada fisicamente na ordem da chave primária (só existe 1 por tabela). \emph{Não-clusterizado} é uma árvore separada ordenada pela coluna indexada que aponta para a chave clusterizada (\emph{Clustered Key}).
    \item \textbf{Para que servem \texttt{INCLUDE}, Composto e Filtrado?}
      \texttt{INCLUDE} anexa colunas somente nas folhas do índice sem inchar os nós internos, zerando \emph{Key Lookup}. Índice composto $(K_1, K_2)$ ordena primeiro por $K_1$ e depois por $K_2$ (igualdade em $K_1$ + faixa em $K_2$ forma leitura contígua). Índice filtrado (\texttt{WHERE Desconto >= 10.00}) indexa apenas o subconjunto de interesse, economizando disco e escrita.
    \item \textbf{O que significa \emph{sargable} (\emph{Search ARGument ABLE})?} É o predicado onde o SGBD consegue fazer busca binária direta na B+ Tree (ex.: \texttt{DataPedido >= '2025-01-01'}). Se aplicarmos função na coluna (\texttt{YEAR(DataPedido) = 2025}) ou curinga à esquerda (\texttt{LIKE N'\%Notebook\%'}), a ordenação do índice deixa de bater e o motor cai para \emph{Scan}.
    \item \textbf{Os 3 joins + \emph{Batch Mode}:} \emph{Nested Loops} (entrada pequena + busca pontual indexada), \emph{Merge Join} (entradas já ordenadas por índice) e \emph{Hash Match} (tabela hash em memória para grandes volumes não ordenados). No \emph{Batch Mode on Rowstore}, operadores como \emph{Hash Match} processam vetores de até $\sim$900 tuplas por vez, economizando muito mais CPU do que o \emph{Row Mode} linha-a-linha.
  \end{itemize}
\end{expBox}

\end{document}
